\scanpage{50}

\[
\Gamma(X'/X)\lhook\joinrel\longrightarrow\Gamma(\unsure{\mathcal X'|\mathfrak X'}/X)\qquad(\text{injectif})
\]

car $X-X_0$ schém. dense dans $X$. Donc $x$ connue quand on connait $z$. Quelle condition à mettre sur $z$~? \deleted{Il faut que} Notons qu'on a

\[
\deleted{\unsure{T_0+z_0}}\qquad T+z\equiv0\quad(2)
\]

car cela signifie $T_0+z_0=0$ i.e. $z_0=T_0$, \inserted{OK}, donc on pourra prendre

\[
x=\frac{T+z}{2},
\]

je dis que a) c'est une section de $P_1$ \unsure{à} i.e. $\psi(x)=1$ (or cela équivaut à $2\psi(x)=2$ i.e. $\underset{\substack{''\\\psi(T+z)}}{\psi(2x)}=2$, Ok),

et b) on a bien $Q(x)=0$ (car ceci équiv. à $4Q(x)=0$ \unsure{i.e.|i.,} $Q(2x)=0$ i.e. $Q(T+z)=0$, donc à $z^2=\delta$ \ldots).

Description de $E$ en termes de $L,\delta,T_0$~: soit \inserted{(loc\textsuperscript{t})} $\tau$ une section de $L$ \inserted{sur $X$} qui remonte $T_0$

\[
\Gamma E=\left\{\underbrace{(\lambda,\unsure{\eta|y|\xi})}_{x}\ \middle|\ \lambda\in\Gamma\underline{\mathscr O},\ x\in\Gamma\left(\underbrace{\frac{\lambda\tau}{2}+L}_{\substack{\text{ne dépend}\\ \text{pas du choix}\\ \text{de }L}}\right)\right\}
\]

\[
\{\lambda=\psi(x),\ \eta=x-\psi(x)\tfrac{T}{2}
\]

\margin{\uline{NB} $\dfrac{\lambda\tau}{2}+L\subset\dfrac12\underline L$}

ou encore

\[
\boxed{\Gamma E\simeq\bigl\{(\lambda,\zeta)\mid\lambda\in\Gamma\underline{\mathscr O},\quad\zeta\in\Gamma(\underbrace{\lambda\tau+2L}_{\overset{\mathrm{déf}}=\ L_\lambda\text{ sous-}\deleted{\unsure{\text{Module}}}\text{torseur de }\underline L})\subset\Gamma(L)}
\]

* i.e. section de $L$ qui remonte la section $\lambda T_0$ de $L_0$ i.e. $\boxed{\zeta_0=\lambda_0T_0}$

\[
\boxed{\left\{\begin{aligned}\lambda&=\psi(x)\\\zeta&=2x-\psi(x)T\end{aligned}\right.}
\]

% slanted
\margin{* via l'op. de $\underline L$ dessus~: $U_\zeta x=x+\underline{2}\zeta$}

Relation entre $(\lambda,\eta)$ et $(\lambda\unsure{,}\zeta)$~: $\zeta=2\eta$

\margin{\deleted{\uline{NB} \unsure{On a} $\boxed{\left\{\begin{aligned}\lambda&=\psi(x)\\\zeta&=2x-\psi(x)T\end{aligned}\right.}$ pour $x=(\lambda,\zeta)$}}

\[
0\longrightarrow\underline L\longrightarrow E\xrightarrow{\psi}\underline{\mathscr O}\longrightarrow0\qquad\left\{\begin{aligned}\psi(\lambda,\zeta)&=\lambda\\i(\unsure{\eta|y|\xi})&=(0,2\eta)\end{aligned}\right.
\]

$T$ correspond à la section $(2,\deleted{\unsure{\zeta}}0)$ de $\Gamma P_2\simeq\Gamma(\underbrace{2\tau+2L}_{2L})$

% archive sheet number: 49
